\documentclass[11pt,a4paper]{ctexart}
\usepackage[a4paper,margin=1.75cm,headheight=14pt]{geometry}
\usepackage{amsmath,amssymb,mathtools}
\usepackage{booktabs,tabularx,longtable,array,multirow}
\usepackage{enumitem}
\usepackage{tikz}
\usetikzlibrary{arrows.meta,positioning,shapes.geometric,fit,calc}
\usepackage[most]{tcolorbox}
\usepackage{xcolor}
\usepackage{fancyhdr}
\usepackage{hyperref}
\usepackage{microtype}
\usepackage{fvextra}
\usepackage{lastpage}
\setlength{\parindent}{2em}
\setlength{\parskip}{0.35em}
\setlist[itemize]{itemsep=0.18em,topsep=0.25em}
\setlist[enumerate]{itemsep=0.18em,topsep=0.25em}
\hypersetup{colorlinks=true,linkcolor=blue!45!black,urlcolor=blue!50!black}

\definecolor{mainblue}{RGB}{45,86,140}
\definecolor{deepblue}{RGB}{30,62,102}
\definecolor{softblue}{RGB}{238,245,252}
\definecolor{softgreen}{RGB}{238,249,243}
\definecolor{softorange}{RGB}{253,246,233}
\definecolor{softred}{RGB}{253,239,239}
\definecolor{softgray}{RGB}{246,247,249}
\definecolor{softpurple}{RGB}{247,241,252}
\definecolor{deepgreen}{RGB}{35,105,75}
\definecolor{deeporange}{RGB}{165,98,22}
\definecolor{deepred}{RGB}{150,55,55}
\definecolor{deeppurple}{RGB}{105,69,133}

\pagestyle{fancy}
\fancyhf{}
\lhead{408 操作系统专题手册}
\rhead{并发 · 信号量 · PV 操作}
\cfoot{第 \thepage\ 页 / 共 \pageref{LastPage} 页}
\renewcommand{\headrulewidth}{0.4pt}
\setlength{\headheight}{14pt}

\newtcolorbox{corebox}[1][]{colback=softblue,colframe=mainblue,title=#1,fonttitle=\bfseries,boxrule=0.8pt,arc=2mm}
\newtcolorbox{exam}[1][]{colback=softorange,colframe=deeporange,title=#1,fonttitle=\bfseries,boxrule=0.8pt,arc=2mm}
\newtcolorbox{warn}[1][]{colback=softred,colframe=deepred,title=#1,fonttitle=\bfseries,boxrule=0.8pt,arc=2mm}
\newtcolorbox{eng}[1][]{colback=softgreen,colframe=deepgreen,title=#1,fonttitle=\bfseries,boxrule=0.8pt,arc=2mm}
\newtcolorbox{method}[1][]{colback=softgray,colframe=black!45,title=#1,fonttitle=\bfseries,boxrule=0.6pt,arc=2mm}
\newtcolorbox{pattern}[1][]{colback=softpurple,colframe=deeppurple,title=#1,fonttitle=\bfseries,boxrule=0.7pt,arc=2mm}

\newcommand{\kw}[1]{\textbf{#1}}
\newcommand{\code}[1]{\texttt{#1}}
\DefineVerbatimEnvironment{codeblock}{Verbatim}{fontsize=\small,breaklines=true,breakanywhere=true,frame=single,framesep=2mm}

\title{\vspace{-1.2cm}\textbf{操作系统并发专题：信号量与 PV 操作方法论}\\[0.45em]
\large 从执行交错到 408 经典同步题的可执行建模方法}
\author{408 专题手册系列 · OS 并发模块}
\date{}

\begin{document}
\maketitle

\begin{corebox}[定位：PV 不是口诀，而是通用 OS 推理引擎在并发同步下的落地特例]
信号量/PV 绝不是背诵伪代码模板，而是通用操作系统系统推理引擎（OS System Reasoning Engine）在“多线程指令交错”场景下的具体应用：
\[
\boxed{
\begin{aligned}
\text{State (共享状态)}
&\to\text{Interleaving (交错转移)}
\to\text{Failure (失效/坏状态)}\\
&\to\text{Property (Safety / Liveness)}
\to\text{Minimal Mechanism (信号量/PV)}
\to\text{Tradeoff (代价)}
\end{aligned}
}
\]
面对陌生并发题，严格回答解题三问：
\begin{enumerate}[leftmargin=1.5em]
\item \kw{第一问（坏状态）：如果完全不加控制，会出现什么错误交错？}（寻找 Counterexample，如 Lost Update / 缓冲区溢出）。
\item \kw{第二问（性质与不变量）：哪些状态必须始终成立？}
  \begin{itemize}
    \item \kw{Safety（安全性）}：坏事绝不发生（如互斥临界区、$0\le count\le N$）。注意区分\kw{宏观逻辑不变量}（如动作完成后 $empty+full+in\_transition=N$）与\kw{指令中途的快照}；
    \item \kw{Liveness（活性）}：好事最终必发生（如无死锁 No Deadlock、无饥饿 No Starvation）。
  \end{itemize}
\item \kw{第三问（最小机制）：恢复该性质所需的最小约束与 PV 机制是什么？}（只禁止危险执行，不防范过度）。
\end{enumerate}

在本手册中，我们严格区分三个核心视级：
\begin{itemize}[leftmargin=1.5em]
\item \kw{许可计数（Permit Count）}：逻辑世界中的可用资源额度（如 $S=1$ 临界区名额，$S=N$ 空位额度）。
\item \kw{约束映射（Constraint Mapping）}：建模世界中将性质推演翻译为 P/V 逻辑顺序过程。
\item \kw{PV 伪代码（PV Pseudocode）}：实现世界中答题卡上最终写出的 \code{P(S)} / \code{V(S)} 代码。
\end{itemize}
\end{corebox}

\begin{method}[术语预备：并发题先辨认对象]
\begin{itemize}[leftmargin=1.5em]
  \item \kw{并发 (Concurrency)}：多个执行流在同一时间段内交错推进；不要求同一时刻真的由多个 CPU 核同时执行。
  \item \kw{原子操作 (Atomic Operation)}：从其他执行流观察时不可被拆开的状态改变；“一行 C 代码”不自动等于原子操作。
  \item \kw{数据竞争 (Data Race)}：多个执行流并发访问同一共享位置，至少一个访问会写入，且缺少规定顺序或互斥，结果依赖交错顺序。
  \item \kw{信号量 (Semaphore)}：带计数的同步对象；计数表示可用许可数，\code{P} 在无许可时等待并消耗一个许可，\code{V} 归还许可并可能唤醒等待者。
  \item \kw{互斥锁 (Mutex)}：计数上限为 1 的独占锁，保护临界区；它表达“同一时刻最多一个持有者”，不自动表达生产者/消费者条件。
  \item \kw{条件变量 (Condition Variable)}：保存等待关系而不是业务条件本身；线程被唤醒后必须重新检查受锁保护的共享谓词。
  \item \kw{死锁 (Deadlock)}：一组执行流形成永久等待环，每个成员都持有某资源并等待环中下一成员持有的资源。
  \item \kw{饥饿 (Starvation)}：某执行流长期得不到所需资源或 CPU，即使系统整体仍在推进；它与“系统完全停止”的死锁不同。
\end{itemize}
因此本册先写“共享状态和坏交错”，再选择互斥、条件同步或资源分配机制；不能看到任何等待就直接套死锁结论。
\end{method}

\begin{method}[内容边界]
这里定义信号量的许可语义、PV 建模流程、经典同步模式和正确性检查。进程状态转换由进程与调度专题定义；本专题只使用 Running$\to$Blocked、Blocked$\to$Ready 等状态结果。
\end{method}

\tableofcontents
\newpage

\section{第零章：为什么会有并发问题}

\subsection{并发程序必须对所有可能的交错顺序都正确}

在并发编程中，无法只看静态代码，必须分析动态运行时的指令执行序列：
\begin{itemize}[leftmargin=1.5em]
\item \kw{多线程指令交错（Interleaving）}：并发运行时，CPU 调度器可以在任意时刻切换线程。线程 $T_1$ 和线程 $T_2$ 的指令在时间轴上交叉拼合，称为“指令交错”。
\item \kw{执行轨迹（Execution Trajectory / Trace，记作 $\sigma$）}：把一次运行中真实发生的全局指令顺序记录下来，构成的具体流水账叫“执行轨迹”。
\item \kw{执行轨迹集合（Set of Trajectories，记作 $\Sigma$）}：由于 CPU 调度的随机性，所有可能产生的交错流水账构成集合 $\Sigma = \{\sigma_1, \sigma_2, \ldots\}$。
\item \kw{并发控制的目标}：单次运行成功仅代表集合 $\Sigma$ 中的一条轨迹 $\sigma$ 正确。P/V 操作要限制允许出现的交错顺序，排除会产生数据竞争或死锁的轨迹。
\end{itemize}

设有两个执行流 $T_1,T_2$ 共同操作共享状态 $S$。某一条具体交错轨迹表示为：
\[
\sigma=T_1^1,T_2^1,T_1^2,T_2^2,\ldots
\]
并发程序正确性定义为：对所有允许交错轨迹，系统状态都必须满足安全不变量 $I$：
\[
\boxed{\forall\sigma\in\text{Allowed Interleavings},\quad I(S_\sigma)=\mathrm{true}}
\]

\begingroup
\small
\renewcommand{\arraystretch}{1.25}
\begin{tabularx}{\linewidth}{>{\bfseries}p{0.22\linewidth} X}
\toprule
符号 / 术语 & 严谨定义与直观含义解析 \\
\midrule
$\sigma$ (Sigma) & \kw{某一条具体的执行轨迹}。一次运行中 CPU 按时间顺序执行的指令步骤流水账。 \\
$T_i^j$ & \kw{进程/线程 $T_i$ 的第 $j$ 个原子步骤}。例如 $T_1^1$ 表示线程 1 执行的第一条汇编指令。 \\
$S_\sigma$ & \kw{沿着轨迹 $\sigma$ 执行完毕后的共享状态}。如变量 \code{counter} 的最终值。 \\
$I(S)$ & \kw{系统不变量（Invariant）}。无论何时都绝对不能被破坏的安全业务规则（如 $0 \le count \le N$）。 \\
$\forall\sigma \in \text{Allowed Interleavings}$ & \kw{对允许轨迹集合中的任意一条轨迹}。要求对所有合法交错组合全部成立。 \\
\bottomrule
\end{tabularx}
\endgroup

\begin{corebox}[并发篇第一原则]
并发不是“执行顺序不确定”这么简单，而是：\kw{程序员不能假定调度器会在方便的位置切换；高层语句又可能由多条底层步骤组成。} 因此正确性必须对所有允许交错成立。
\end{corebox}

\subsection{综合题：一条 \code{counter++} 的一生}

程序员看到：
\[
\texttt{counter++}
\]
机器可能执行：
\begin{codeblock}
load counter
add 1
store counter
\end{codeblock}
若初值为 5，两个线程可能发生：
\begin{codeblock}
T1: load 5
T2: load 5
T1: add  -> 6
T2: add  -> 6
T1: store 6
T2: store 6
\end{codeblock}
最终得到 6 而不是 7，称为\kw{丢失更新（Lost Update）}。

这条最小反例依次推出：
\[
\boxed{
\text{Shared State}
\to\text{Interleaving}
\to\text{Race}
\to\text{Atomicity}
\to\text{Critical Section}
\to\text{Mutual Exclusion}
}
\]

\subsection{同步不等于互斥}

并发控制中最基本的两个关系必须分开：

\begingroup
\small
\renewcommand{\arraystretch}{1.25}
\setlength{\tabcolsep}{6pt}
\begin{tabularx}{\linewidth}{>{\bfseries}p{0.24\linewidth} X X}
\toprule
关系类型 & 需要解决的问题 & 典型工具与直观表现\\
\midrule
互斥 (Mutual Exclusion) & 哪些关键指令/状态更新不能彼此交错？ & Mutex、二元信号量 (Initial=1)、临界区\\
条件/顺序同步 (Synchronization) & 什么条件满足后才能继续？谁必须先于谁？ & 计数信号量 (Initial=N/0)、事件信号量、条件变量\\
\bottomrule
\end{tabularx}
\endgroup

\begin{exam}[考试识别]
看到“共享变量、共享缓冲区、文件写入、计数器”等，先找\kw{不允许交错的状态更新}；看到“有产品才能消费、A 完成后 B 才能做、资源数量有限”等，先找\kw{等待条件/前驱关系}。很多 PV 大题是两者叠加。
\end{exam}

\subsection{临界区四段与四条准则：先分 Safety、Progress，再谈效率}

对某一份临界资源，执行流的循环结构应拆成：
\[
\boxed{\text{Entry}\to\text{Critical Section}\to\text{Exit}\to\text{Remainder}}
\]
Entry 负责原子地检查/取得准入，Critical Section 执行不可交错的共享状态事务，Exit 归还资格并通知竞争者，Remainder 则是不需要被这把锁串行化的工作。一份资源可以被不同进程中的多段代码访问，它们必须共享同一准入协议；反过来，不碰同一不变量的代码不应因为“属于同一进程”就被同一把粗锁圈住。

教材四条准则并非同一层：
\begin{enumerate}
\item \kw{忙则等待}：已有执行流在临界区时，其他竞争者不能进入，拥有 Mutual Exclusion / Safety；
\item \kw{空闲让进}：临界区空闲且有请求者时，系统必须能选出一个进入，排除“永远不让任何人进”的伪安全方案；
\item \kw{有限等待}：某个持续请求者不能被无限次越过，要求有界等待/无饥饿所需的公平假设；
\item \kw{让权等待}：暂时不能进入者应阻塞让出 CPU，而非无限 busy-wait；它主要是资源效率要求。
\end{enumerate}
判陌生方案先攻击第 1 条；Safety 已破就无需再用“看起来公平”挽救。满足前三条但忙等的 Peterson/TAS 在经典模型下仍可正确，只是未满足第 4 条。第 4 条需要 block/wakeup 与调度支持，纯用户态普通变量读写本身无法把当前实体移入内核等待队列。

\subsection{软件互斥四步修补链：每一次修补都暴露下一种失败}

两进程经典算法可以按失败反例生成，而不是背四段代码：
\begin{enumerate}
\item \kw{单标志 turn}：互斥成立，却强制严格交替；另一方不再请求时，本方仍不能连续进入，违反空闲让进；
\item \kw{双标志先检查后举手}：消除强制交替，但两方可先同时读到对方未举手，再同时置位进入，违反忙则等待；
\item \kw{双标志先举手后检查}：修复同时进入，却可能两方同时举手后都等待对方撤旗，临界区空闲而无人能进；
\item \kw{Peterson = flag + turn}：flag 表达意愿，turn 在双方同时竞争时打破对称；后写“让对方先”的进程自我等待，另一方推进后撤旗。
\end{enumerate}
统一反例控制器是把每方的“检查/声明”拆成两拍，尝试交错
\[
a_1\to b_1\to a_2\to b_2.
\]
先检查法由它构造双入；后检查法由它构造双等；Peterson 由唯一的 turn 值裁决一方。该正确性还依赖所采用的共享内存模型能提供算法要求的原子读写与顺序；现代弱内存架构和编译器优化下需要相应原子/内存序语义，不能把源码普通变量版本直接当跨平台锁实现。

\subsection{互斥的底座机制：关中断 vs 硬件原子指令}

为了实现互斥（Critical Section），硬件与系统提供了从单核到多核的底座机制演进：

\begin{enumerate}
  \item \kw{关中断 (Disable Interrupt)}：单 CPU 运行环境下，CPU 通过关闭本地中断（cli），防止时钟中断触发调度切换，确保临界区代码不被同一 CPU 上的其他线程抢占。
  \item \kw{硬件原子指令 (TestAndSet / CAS / LL-SC)}：现代多核 (SMP) 架构下，多 CPU 共享主存。关闭本地 CPU 的中断**无法阻止其他 CPU 核心访问共享内存**！必须依赖 CPU 硬件层面的原子读改写（RMW）指令与缓存一致性/总线锁。
\end{enumerate}

必须确立核心判定边界：
\[
\boxed{
\text{Disable Local Interrupt}
\not\Rightarrow
\text{Multiprocessor Global Mutual Exclusion}
}
\]

\begingroup
\small
\renewcommand{\arraystretch}{1.25}
\setlength{\tabcolsep}{4pt}
\begin{tabularx}{\linewidth}{>{\bfseries}p{0.20\linewidth} X X}
\toprule
机制 / 方法 & 核心原理与适用场景 & 缺陷与代价\\
\midrule
Disable Interrupt & 关本地 CPU 中断，防调度抢占；适用于单核 CPU 极短内核操作 & 多核 (SMP) 无效；滥用关中断可能导致系统响应丧失与实时性破坏\\
Software Peterson & 纯软件标志位 (flag + turn) 实现两线程互斥 & 存在 Busy-waiting（忙等）；依赖内存顺序，难以推广到 $N$ 线程\\
TestAndSet (TAS) & 硬件指令原子性地“读旧值并写新值 1” & 存在 Busy-waiting 消耗 CPU 算力（Spinlock 自旋锁特征）\\
Compare-And-Swap & 比较当前值是否等于期望值，相等则原子更新 & 存在 ABA 问题（需版本号解决）；自旋开销受高并发争用影响\\
\bottomrule
\end{tabularx}
\endgroup

\subsection{Spin 还是 Sleep：比较等待预算，而不是按锁名字站队}

自旋锁拿不到时保持 runnable 并反复检查；睡眠 mutex 把等待者挂入锁等待队列，释放时再 wake 到 Ready。若预计剩余持锁时间为 $T_h$，阻塞、调度、恢复和局部性重热的总成本为 $C_b$，最小判断是：
\[
\boxed{T_h<C_b\Rightarrow\text{短时自旋可能更便宜};\qquad
T_h\gg C_b\Rightarrow\text{睡眠通常更合算}.}
\]
单 CPU 上若自旋者抢占了唯一能释放锁的持有者，自旋不会制造进度；多 CPU 上持有者可能在别的核继续运行，自旋才有收益机会。中断/原子上下文不能睡眠时也可能只能用自旋，但这不是“自旋免费”。

朴素 TAS 每轮都写锁变量，会反复争夺缓存行。Test-and-Test-and-Set 先只读观察，只有看见空闲才执行一次 RMW：
\[
\text{read-only spin}\to\text{appears free}\to\text{TAS attempt}.
\]
它减少锁仍被持有期间的写失效风暴，却不能消除同时看到空闲后的争抢。锁粒度越细，互不相关操作的并行度越高，但一次操作可能取得更多锁，排序、回滚与死锁证明成本随之增加。

\subsection{优先级反转与两类锁协议：先画三方依赖，再比较触发时机}

优先级反转的最小链条需要三方：
\[
H\to\text{Lock held by }L,\qquad M\succ L,\qquad M\text{ does not need that lock}.
\]
$H$ 因锁直接等 $L$，$M$ 又抢占 $L$，使 $H$ 被无直接资源关系的 $M$ 间接延迟。只有 $H,L$ 两方时，高优先级等待仍是锁阻塞，但缺少 $M$ 放大的“反转”链。

\begin{itemize}
\item \kw{Priority Inheritance}：$H$ 真正阻塞后，持锁者 $L$ 临时继承等待者中的最高有效优先级，释放相关锁后按仍存在的依赖恢复；若 $L$ 又等待 $X$ 持有的锁，提升必须沿 wait-for chain 传递。它动态、按需，但实现必须维护嵌套锁和多等待者关系；它本身不禁止循环取锁，所以不能替代 deadlock prevention。
\item \kw{Priority Ceiling Protocol Family}：预先为锁记录所有潜在使用者的最高优先级，并按具体 ceiling 协议限制取得锁或在取得时立即提升。静态可分析的实时任务集可据此限制阻塞链，并在相应协议假设下给出“一次阻塞”等上界及防止特定锁死锁；代价是需要预知访问集合，且可能提前提升本来不会被高优先级竞争的持锁者。
\end{itemize}
“拿锁即升到 ceiling”描述的是常见 immediate-ceiling 变体，不代表所有 ceiling 协议的逐步规则。考试若给出具体协议，以它对系统 ceiling、准入和恢复的定义为准。

\section{信号量到底是什么：把并发约束映射为许可计数}

\subsection{最稳定的理解：Permit Counter + Waiting Queue}

信号量可以理解为：
\[
\boxed{\text{Semaphore}=\text{Permit Counter}+\text{Atomic Wait/Wake Mechanism}}
\]
它不是“神秘整数”，而是把某种现实业务规则映射为具体可消费的许可数量（Permit Count）：

\begin{itemize}
\item $S=1$：当前有 1 个进入资格，可用于单资源互斥；
\item $S=N$：当前有 $N$ 个同类资源或容量许可；
\item $S=0$：某个必须等待的事件尚未发生，或当前没有可消费资源。
\end{itemize}

\begin{center}

\begingroup
\normalsize
\begin{tikzpicture}[>=Stealth, node distance=1.6cm and 3.0cm]
\tikzset{
  box/.style={draw=mainblue, ultra thick, fill=softblue, rounded corners=3mm, align=center, minimum width=3.2cm, minimum height=1.1cm, font=\bfseries},
  act/.style={draw=deeporange, ultra thick, fill=softorange, rounded corners=3mm, align=center, minimum width=3.0cm, minimum height=1.0cm, font=\small\bfseries}
}
\node[box] (sem) {信号量实体 $S$\\[3pt]\small\normalfont 许可计数器 Value};
\node[box, right=3.5cm of sem] (queue) {阻塞等待队列\\[3pt]\small\normalfont Wait Queue (PCBs)};

\node[act, above left=1.2cm and 1.2cm of sem] (p_op) {\code{P(S) / wait()}};
\node[act, below left=1.2cm and 1.2cm of sem] (v_op) {\code{V(S) / signal()}};

\draw[->, ultra thick, deepblue] (p_op) -- node[above=2pt, sloped, font=\small\bfseries] {$S > 0 \implies S--$} (sem);
\draw[->, ultra thick, deepred] (p_op) -| node[above=4pt, font=\small\bfseries] {$S \le 0 \implies \text{当前线程入队阻塞 (Blocked)}$} (queue);

\draw[->, ultra thick, deepblue] (v_op) -- node[below=2pt, sloped, font=\small\bfseries] {$S++$} (sem);
\draw[->, ultra thick, deepgreen] (queue) |- node[below=4pt, font=\small\bfseries] {队列非空 $\implies$ 唤醒队首线程 (Ready)} (v_op);
\end{tikzpicture}
\endgroup
\end{center}

\begin{corebox}[初值不是背出来的]
问：\kw{系统初始时到底已经拥有多少份可以立即被消费的事实/资源/空位/许可？} 这个数量就是计数信号量初值的来源。
\end{corebox}

\subsection{P/V 的两种常见教材记法}

不同教材会采用不同内部表示。做题时首先服从题目定义。

\paragraph{记法 A：可出现负值的经典教学模型}
\begin{align*}
P(S):&\quad S\leftarrow S-1;\quad S<0\Rightarrow\text{当前进程阻塞入队}\,,\\
V(S):&\quad S\leftarrow S+1;\quad S\le 0\Rightarrow\text{从等待队列唤醒一个进程}\,.
\end{align*}
在这种模型中，负值的绝对值常可解释为等待者数量。

\paragraph{记法 B：计数保持非负的抽象模型}
\begin{itemize}
\item \code{wait/P}：若有许可则原子地消费一个；若无许可则阻塞；
\item \code{post/V}：增加一个许可；若有等待者则允许其中一个继续竞争运行。
\end{itemize}

两者语义一致：\kw{P 可能阻塞，V 可能唤醒，且 P/V 自身必须原子。}

\subsection{从整型到记录型：资源不足时进入等待队列}

经典整型信号量写作
\begin{codeblock}
wait(S)   { while (S <= 0) ; S--; }
signal(S) { S++; }
\end{codeblock}
即使把“检查并减一”视为一个可信原子操作，等待者仍在 while 中占用 CPU。记录型信号量增加等待队列，并把 P/V 接入 OS 的 block/wakeup：
\begin{codeblock}
P(S) {
    S.value--;
    if (S.value < 0) block(S.queue);
}
V(S) {
    S.value++;
    if (S.value <= 0) wakeup_one(S.queue);
}
\end{codeblock}
在这套负值教材记法中，$S.value<0$ 时 $|S.value|$ 常表示等待者数量。让权等待的关键不是“结构体里多了一根链表指针”，而是\kw{把“修改计数 + 加入等待队列 + 阻塞”做成无 lost wakeup 的原子状态事务}；V 则把一份许可交给一个等待者并使其 Ready。具体实现也可以保持非负计数并单独记录 waiter，语义不依赖负值编码。

\subsection{AND 型信号量：把多资源检查与取得合成一次事务}

若进程逐个执行 $P(S_1),P(S_2)$，可能在已持有 $S_1$ 时阻塞等待 $S_2$，与反序取得者形成 Hold-and-Wait 环。AND 型操作要求：
\[
\boxed{
\operatorname{Swait}(S_1,\ldots,S_n):
\left(\bigwedge_i S_i\ge1\right)
\Rightarrow
\forall i,\ S_i\leftarrow S_i-1
}
\]
若任一项不足，则全部计数保持不变并等待；释放以同一多对象协议归还。它消除“只拿到子集”的可观察中间态，从而在这组资源上破坏请求并保持。但实现本身必须为多个信号量建立统一原子化/锁序，不能用普通 P 的循环拼接后仍宣称 AND 语义。

\subsection{信号量集：把“门槛”与“消耗量”分开}

教材信号量集把一次操作写成三元组 $(S_i,t_i,d_i)$：
\[
\operatorname{Swait}(S_i,t_i,d_i):
\quad
\bigwedge_i(S_i\ge t_i)
\Rightarrow
S_i\leftarrow S_i-d_i.
\]
$t_i$ 是允许事务发生的测试下限，$d_i$ 是真正消耗量；$d_i=0$ 因而能表达“只测试，不占用”。典型退化关系是：
\begin{itemize}
\item $\operatorname{Swait}(S,d,d)$：一次申请 $d$ 份；
\item $\operatorname{Swait}(S,1,1)$：普通单份申请；
\item $\operatorname{Swait}(S,1,0)$：把 $S$ 当可控闸门，只检查开放状态。
\end{itemize}
这些是特定教材 API 的能力定义，不代表普通 semaphore wait 可以随意“只测不减”。做题先确认题目是否真的提供 Swait/Ssignal，再使用该表达力。

\subsection{P/V 为什么必须是原子操作}

如果 \code{P(S)} 自己可以被拆成“读 S、判断、减一、入队”等可交错步骤，那么两个执行流可能同时认为自己拿到了同一份许可。信号量的正确性因此必须最终依赖更底层的原子机制，如 test-and-set、compare-and-swap、fetch-and-add 或内核的受保护临界区。

\begin{pattern}[抽象层次]
\[
\boxed{
\text{硬件原子能力}
\to\text{锁/阻塞原语}
\to\text{信号量/条件变量}
\to\text{并发算法}
}
\]
408 通常从信号量层开始做题，但理解下面一层能解释“为什么 P/V 本身可以被信任”。
\end{pattern}

\subsection{二元信号量与 Mutex：考试等效，语义并非完全相同}

408 里常说“二元信号量可实现互斥”，这是正确的考试模型。但严格设计上：
\begin{itemize}
\item mutex 强调\kw{所有权（ownership）}：谁加锁，通常由谁解锁；
\item semaphore 强调\kw{许可数量}：一个执行流 P，另一个执行流 V 完全可以有合理语义，例如事件同步。
\end{itemize}
因此应记：
\[
\boxed{\text{Binary Semaphore can implement mutual exclusion, but is not conceptually identical to a mutex.}}
\]

\section{最小调用接口：从共享状态进入许可模型}

当题目出现多个执行流、共享状态、等待条件或资源计数时，先回到本册确认三件事：\kw{谁在等待什么事实、谁能制造这个事实、哪些共享更新必须保持原子不变量}。本册负责解释 semaphore/PV 为什么能够表达这些约束，以及等待依赖如何产生死锁；完整的落笔流程、P/V 排序检查、最少信号量与作答控制统一进入同目录训练文档 \texttt{PV同步题.md}。

若问题已经从同步约束转为“当前状态是否安全、请求能否批准、是否已经死锁或怎样恢复”，转入 \texttt{死锁判断与Banker.md}；Canonical 继续拥有 Safety、Banker、Detection 与 Recovery 的机制语义。

\section{模式一：生产者--消费者——容量计数 + 互斥更新}

\subsection{问题模型与不变量}

容量为 $N$ 的有界缓冲区：
\[
0\le count\le N
\]
生产者继续条件：
\[
count<N
\]
消费者继续条件：
\[
count>0
\]
同时，插入/删除和缓冲区元数据更新不能被多个执行流破坏。

\subsection{信号量设计}

\begingroup
\small
\renewcommand{\arraystretch}{1.25}
\setlength{\tabcolsep}{6pt}
\begin{tabularx}{\linewidth}{>{\bfseries}p{0.2\linewidth} p{0.18\linewidth} X}
\toprule
信号量对象 & 初值 (Initial) & 现实含义与许可类型\\
\midrule
\code{empty} & $N$ & 可立即消费的空位许可数量（同步信号量）\\
\code{full} & $0$ & 可立即消费的产品许可数量（同步信号量）\\
\code{mutex} & $1$ & 进入缓冲区结构更新临界区的唯一许可（互斥信号量）\\
\bottomrule
\end{tabularx}
\endgroup

关键守恒关系可辅助检查：
\[
empty+full=N
\]
（在经典抽象模型、忽略正在 P/V 过渡的瞬间进行逻辑理解。）

\subsection{标准模板}

\begin{codeblock}
semaphore mutex = 1;
semaphore empty = N;
semaphore full  = 0;

producer() {
    while (true) {
        produce_item();
        P(empty);
        P(mutex);
        put_item();
        V(mutex);
        V(full);
    }
}

consumer() {
    while (true) {
        P(full);
        P(mutex);
        get_item();
        V(mutex);
        V(empty);
        consume_item();
    }
}
\end{codeblock}

\subsection{为什么生产和消费要放在锁外}

\code{produce\_item()} 与 \code{consume\_item()} 若只使用私有数据，不应放入缓冲区 mutex 临界区。否则虽然 Safety 仍可正确，却会把长时间工作串行化，降低并发度。

\subsection{经典死锁轨迹}

缓冲区已满，若生产者先：
\begin{codeblock}
P(mutex);   // 成功
P(empty);   // empty=0 -> Blocked
\end{codeblock}
则生产者睡眠时仍持有 mutex。消费者即使 Ready/Running，也会阻塞在 \code{P(mutex)}，无法取出产品并执行 \code{V(empty)}。形成：

\begin{center}

\begingroup
\normalsize
\begin{tikzpicture}[>=Stealth, node distance=1.8cm and 3.0cm]
\tikzset{
  proc/.style={draw=deepred, ultra thick, fill=softred, rounded corners=3mm, align=center, minimum width=3.2cm, minimum height=1.1cm, font=\bfseries},
  res/.style={draw=mainblue, ultra thick, fill=softblue, rounded corners=3mm, align=center, minimum width=3.0cm, minimum height=1.1cm, font=\small\bfseries}
}
\node[proc] (p) {生产者 Producer\\[3pt]\small\normalfont 持有 \code{mutex} 锁};
\node[res, right=3.5cm of p] (e) {等待条件 \code{empty=0}\\[3pt]\small\normalfont 需 Consumer 执行 V};
\node[proc, below=1.8cm of e] (c) {消费者 Consumer\\[3pt]\small\normalfont 阻塞在 \code{mutex} 锁};
\node[res, below=1.8cm of p] (m) {互斥锁 \code{mutex}\\[3pt]\small\normalfont 被 Producer 独占};

\draw[->, ultra thick, deepred] (p) -- node[above=2pt, font=\small\bfseries] {持有锁并等待条件} (e);
\draw[->, ultra thick, deepred] (e) -- node[right=2pt, font=\small\bfseries] {只能由 Consumer 触发 V} (c);
\draw[->, ultra thick, deepred] (c) -- node[below=2pt, font=\small\bfseries] {欲取锁而无法进入} (m);
\draw[->, ultra thick, deepred] (m) -- node[left=2pt, font=\small\bfseries] {锁被 Producer 占用} (p);
\end{tikzpicture}
\endgroup
\end{center}

这会形成等待环。

\subsection{典型变式}

\begin{itemize}
\item 单/多生产者、单/多消费者：基本模板不变；
\item 消费一个 A 和一个 B：分别建立 \code{fullA}、\code{fullB}，消费者要取得两类许可；
\item 多类资源同时获取：开始接近哲学家/多资源分配问题，需要检查循环等待和持有后等待；
\item 批量生产/消费 $k$ 个：不能只机械连续执行 $k$ 次 P；要明确题目是否要求“整批原子获得”，以及容量/公平性要求。
\end{itemize}

\begin{exam}[容量为 1 时为什么有时可省 mutex，容量大于 1 又为什么不行]
单槽“盘子”若所有生产/消费都先取得同一个初值为 1 的 slot/state 许可，许可本身已经把对盘子的结构更新串行化，额外 mutex 可能冗余。但 $N>1$ 时，\code{empty} 允许多名生产者同时通过，\code{full} 也允许多名消费者同时通过；它们只约束容量，不保护共享 \code{in/out} 与槽位元数据，因此仍需 mutex 或更细粒度的并发队列协议。
\end{exam}

\begin{method}[多类产品：信号量正确不代表数据结构已经能定位对象]
多类消费者通常共享一个 \code{empty=N}，每类产品再有 \code{fullA/fullB/...=0}；消费者先取得对应类型许可，再互斥修改容器。但若底层只是严格 FIFO 的单一 \code{out} 指针，拿到 \code{fullA} 不保证队头就是 A。实现必须另设分类队列、类型索引或可搜索槽位。这里揭示两层不变量：
\[
\boxed{\text{Availability Accounting}\neq\text{Object Location/Removal}}
\]
PV 许可计数正确，只证明“某类对象存在”；还必须证明容器操作能找到并原子移除同一个对象。
\end{method}

\section{模式二：读者--写者——群体准入 + 元数据保护}

\subsection{约束结构}

约束：
\begin{itemize}
\item 读--读允许；
\item 读--写互斥；
\item 写--写互斥。
\end{itemize}
它不是普通“一把锁保护所有人”，而是\kw{群体准入（group admission）}：多个读者构成一个共享访问群体，第一个读者负责关闭写者入口，最后一个读者负责重新开放。

\subsection{读者优先模板}

\begin{codeblock}
semaphore rw_mutex = 1;  // 数据区读写锁
semaphore r_mutex  = 1;  // 保护 readcount
int readcount = 0;

reader() {
    P(r_mutex);
    if (readcount == 0)
        P(rw_mutex);      // 第一个读者封锁写者
    readcount++;
    V(r_mutex);

    read_data();

    P(r_mutex);
    readcount--;
    if (readcount == 0)
        V(rw_mutex);      // 最后一个读者释放写者
    V(r_mutex);
}

writer() {
    P(rw_mutex);
    write_data();
    V(rw_mutex);
}
\end{codeblock}

\subsection{为什么 \code{readcount} 本身必须保护}

\code{readcount++}、判断“是否第一个读者”、\code{readcount--}、判断“是否最后一个读者”共同决定群体状态。如果它们交错，就可能让多个读者都误以为自己是第一/最后一个，破坏准入协议。

\begin{corebox}[读者--写者问题要检查的三件事]
它训练的不是“多加一把锁”，而是：\kw{共享数据和管理共享数据的元数据都可能是临界状态。} 元数据（如计数器、索引、状态位）常常比业务数据本身更容易被忽略。
\end{corebox}

\subsection{公平性是独立策略维度}

上述方案是读者优先：只要读者不断到来，写者可能长期无法获得 \code{rw\_mutex}，即\kw{写者饥饿}。

可以再增加“入口闸门”和写者计数实现写者优先：

\begin{codeblock}
int read_count = 0;
int write_count = 0;
semaphore r_mutex = 1;
semaphore w_mutex = 1;
semaphore rw_mutex = 1;
semaphore read_try = 1;

writer() {
    P(w_mutex);
    write_count++;
    if (write_count == 1) P(read_try);
    V(w_mutex);

    P(rw_mutex);
    write_database();
    V(rw_mutex);

    P(w_mutex);
    write_count--;
    if (write_count == 0) V(read_try);
    V(w_mutex);
}

reader() {
    P(read_try);
    P(r_mutex);
    read_count++;
    if (read_count == 1) P(rw_mutex);
    V(r_mutex);
    V(read_try);

    read_database();

    P(r_mutex);
    read_count--;
    if (read_count == 0) V(rw_mutex);
    V(r_mutex);
}
\end{codeblock}

\begin{exam}[考试层次]
408 常见题主要要求读者优先、写者优先或给定策略下填 PV。解题时先确定题目要求的\kw{公平策略}，再决定是否需要额外闸门；不要把某一版模板当成唯一“读者--写者答案”。
\end{exam}

\subsection{读写公平：让两类请求共用到达闸门}

若读者与写者先在同一 \code{service\_queue} 上按题设队列顺序通过，后到读者不能越过已排队写者并接到现有读者群后面：
\begin{codeblock}
semaphore service_queue = 1;
semaphore rw_mutex = 1;
semaphore r_mutex = 1;
int read_count = 0;

reader() {
    P(service_queue);
    P(r_mutex);
    if (++read_count == 1) P(rw_mutex);
    V(service_queue);
    V(r_mutex);
    read_data();
    P(r_mutex);
    if (--read_count == 0) V(rw_mutex);
    V(r_mutex);
}

writer() {
    P(service_queue);
    P(rw_mutex);
    V(service_queue);
    write_data();
    V(rw_mutex);
}
\end{codeblock}
它把“谁先形成准入承诺”交给公共闸门，而读者通过后尽快释放闸门，仍允许同一批读者并发。无饥饿结论还依赖 \code{service\_queue} 的等待队列具有公平/有界越过语义；semaphore API 若任意选择唤醒者，只看伪代码不能证明严格 FIFO。

\section{模式三：理发师——有限准入 + 双向会合}

\subsection{模型抽象}

理发店有 $M$ 个等待座位。顾客到达：有空位则等待，无空位则离开；理发师无顾客时睡眠。

这个题可以理解为：
\begin{itemize}
\item 顾客产生服务请求；
\item 理发师消费服务请求；
\item 等待椅是有限容量；
\item 顾客与理发师还需完成一次\kw{rendezvous（会合）}。
\end{itemize}

\subsection{经典信号量}

\begingroup
\small
\renewcommand{\arraystretch}{1.25}
\setlength{\tabcolsep}{6pt}
\begin{tabularx}{\linewidth}{>{\bfseries}p{0.28\linewidth} p{0.15\linewidth} X}
\toprule
信号量 / 共享变量 & 初值 (Initial) & 含义与作用\\
\midrule
\code{customers} & 0 & 正在等待服务的顾客数量（同步信号量）\\
\code{barbers} & 0 & 已准备好接待顾客的理发师许可（同步信号量 / 双向会合）\\
\code{mutex} & 1 & 保护 \code{waiting\_customers} 计数器的互斥信号量\\
\code{waiting\_customers} & 0 & 当前在等待椅子上坐着的顾客人数（共享整型变量）\\
\bottomrule
\end{tabularx}
\endgroup

\begin{codeblock}
barber() {
    while (true) {
        P(customers);
        P(mutex);
        waiting_customers--;
        V(barbers);
        V(mutex);
        cut_hair();
    }
}

customer() {
    P(mutex);
    if (waiting_customers < CHAIRS) {
        waiting_customers++;
        V(customers);
        V(mutex);
        P(barbers);
        get_haircut();
    } else {
        V(mutex);
        leave_shop();
    }
}
\end{codeblock}

\subsection{它为什么比生产者--消费者多一层}

生产者--消费者只需确保“产品存在”；理发师问题还要让一个具体顾客知道“已有理发师可以服务”，因此多了 \code{barbers} 这一会合信号。题目若增加多名理发师、不同服务等级或排队规则，就会进一步引入身份匹配与公平性。

\section{模式四：哲学家进餐——多资源获取 + 循环等待}

\subsection{把哲学家问题画成资源依赖图}

最朴素方案：每人先左后右。若所有哲学家都拿到左筷子：
\[
P_0\to C_1\to P_1\to C_2\to\cdots\to P_{N-1}\to C_0\to P_0
\]
形成循环等待。

\subsection{常见死锁规避思路}

\begin{enumerate}
\item \kw{限制准入}：最多允许 $N-1$ 人同时尝试拿筷子；
\item \kw{资源全序}：所有人按统一编号顺序获取资源，破坏循环等待；
\item \kw{原子化资格判定}：在受保护状态中判断左右邻居，只有同时满足才进入 EATING；否则阻塞等待通知。
\end{enumerate}

\begin{method}[四类经典解法究竟破坏了哪个中间状态]
\begin{itemize}
\item 准入信号量初值 $N-1$：最多 $N-1$ 人进入取资源阶段，保证至少有一人能取得第二份资源，阻断全环闭合；
\item AND/Swait 同时取得左右两份：要么全拿、要么一份不持有，直接消除“持一等一”；
\item 奇偶/全序取法：改变局部申请方向，使全局 wait edge 不可能按同一方向闭环；
\item 用 mutex 包住“两次取得”：把取两份资格串成原子阶段，同样消除多个哲学家各持一份的状态，但某执行流可能在持有该 mutex 时等待筷子，降低并发且必须证明不会把新的锁依赖环带进来。
\end{itemize}
前两类分别从准入数量与 Hold-and-Wait 入手，第三类破 Circular Wait，第四类用粗粒度串行换简单性。它们都可以无死锁，但不自动证明无饥饿；最大同时进餐人数主要受资源拓扑约束，也不是评价公平性的充分指标。
\end{method}

\subsection{状态数组 + 私有信号量的阻塞范式}

\begin{codeblock}
semaphore mutex = 1;
semaphore self[N] = {0};
state[N] = {THINKING};

test(i) {
    left  = (i + N - 1) % N;
    right = (i + 1) % N;
    if (state[i] == HUNGRY &&
        state[left]  != EATING &&
        state[right] != EATING) {
        state[i] = EATING;
        V(self[i]);
    }
}

take_forks(i) {
    P(mutex);
    state[i] = HUNGRY;
    test(i);
    V(mutex);
    P(self[i]);
}

put_forks(i) {
    P(mutex);
    state[i] = THINKING;
    test(left(i));
    test(right(i));
    V(mutex);
}
\end{codeblock}

这里 \code{mutex} 保护的是\kw{状态判定过程}，\code{self[i]} 表达“哲学家 $i$ 何时获得进餐资格”。

\begin{warn}[不要过度声称]
这种状态判定方案可以设计为\kw{无死锁}，但“绝对无饥饿”还依赖唤醒/调度公平性和具体判定策略。考试若只问死锁避免，不必额外假设强公平性。
\end{warn}

\section{模式五：吸烟者——条件分派与精确通知}

\subsection{核心不是“材料”，而是多路条件}

供应者每轮制造不同条件，只有匹配的一个吸烟者应继续：
\[
Condition_0\to Smoker_0,\quad
Condition_1\to Smoker_1,\quad
Condition_2\to Smoker_2.
\]
最直接的 PV 思路是：\kw{为每种可区分事件设置独立同步信号量}。

\subsection{简洁 408 模板}

设 \code{agent=1} 表示桌面空、允许供应下一组；\code{offer[3]=0} 表示三种匹配事件尚未发生。

\begin{codeblock}
provider() {
    while (true) {
        P(agent);
        choose k in {0,1,2};
        put_materials_for(k);
        V(offer[k]);
    }
}

smoker(k) {
    while (true) {
        P(offer[k]);
        take_materials();
        make_and_smoke();
        V(agent);
    }
}
\end{codeblock}

这体现：
\[
\boxed{\text{复杂条件同步}\to\text{把不同条件编码为不同等待队列/信号量}}
\]

\begin{eng}[关于“代理人模式”]
增加 Agent/检查者可以把匹配分发移入中间层，但不是 408 必需模板。若多个代理共同消费同一检查信号，必须防止错误消费事件或形成无效唤醒循环；考试中优先为每类条件设置独立信号量。
\end{eng}

\section{模式六：前驱图——把 Happens-Before 边编码成事件信号量}

\subsection{图模型}

若有有向边：
\[
A\to C
\]
表示 $A$ 完成是 $C$ 开始的前置事实。若两操作位于不同执行流，就需要显式同步；若在同一线程中且程序顺序已保证，则通常无需额外信号量。

\begin{corebox}[一条跨线程前驱边 = 一个可消费事件]
为跨线程边 $A\to C$ 定义初值为 0 的信号量 $S_{AC}$：
\[
A;\ V(S_{AC})
\qquad\qquad
P(S_{AC});\ C
\]
初值为 0 正是因为“$A$ 已完成”这一事实在系统初始时尚不存在。
\end{corebox}

\subsection{2022 风格示例}

依赖：
\[
A\to C,\ B\to C,\ C\to D,\ C\to E,\ E\to F
\]
线程：
\[
T_1=\{A,E,F\},\qquad T_2=\{B,C,D\}.
\]
只保留跨线程边：
\[
A\to C,\qquad C\to E.
\]

\begin{center}

\begingroup
\normalsize
\begin{tikzpicture}[>=Stealth, node distance=1.8cm and 2.5cm]
\tikzset{
  t1node/.style={draw=mainblue, ultra thick, fill=softblue, rounded corners=2.5mm, align=center, minimum width=2.2cm, minimum height=0.95cm, font=\normalsize\bfseries},
  t2node/.style={draw=deeporange, ultra thick, fill=softorange, rounded corners=2.5mm, align=center, minimum width=2.2cm, minimum height=0.95cm, font=\normalsize\bfseries}
}
\node[t1node] (a) {$A$};
\node[t1node, right=of a] (e) {$E$};
\node[t1node, right=of e] (f) {$F$};

\node[t2node, below=2.0cm of a] (b) {$B$};
\node[t2node, right=of b] (c) {$C$};
\node[t2node, right=of c] (d) {$D$};

\draw[->, ultra thick, deepblue] (a) -- node[above=3pt, font=\small] {线程 $T_1$ 程序顺序} (e);
\draw[->, ultra thick, deepblue] (e) -- node[above=3pt, font=\small] {线程 $T_1$ 程序顺序} (f);

\draw[->, ultra thick, deeporange] (b) -- node[below=3pt, font=\small] {线程 $T_2$ 程序顺序} (c);
\draw[->, ultra thick, deeporange] (c) -- node[below=3pt, font=\small] {线程 $T_2$ 程序顺序} (d);

\draw[->, ultra thick, deepred] (a) -- node[left=4pt, font=\small\bfseries] {\code{V(ac)} $\rightarrow$ \code{P(ac)}} (c);
\draw[->, ultra thick, deepgreen] (c) -- node[left=4pt, font=\small\bfseries] {\code{V(ce)} $\rightarrow$ \code{P(ce)}} (e);
\end{tikzpicture}
\endgroup
\end{center}

于是：
\begin{codeblock}
semaphore ac = 0, ce = 0;

T1() {
    A();
    V(ac);
    P(ce);
    E();
    F();
}

T2() {
    B();
    P(ac);
    C();
    V(ce);
    D();
}
\end{codeblock}

\subsection{新增模式：Barrier 与 fork/join}

\paragraph{Fork/Join}
线程 A 创建工作线程 B，A 必须等待 B 完成：这就是一次性事件。
\[
S_{done}=0,\quad B:\ V(S_{done}),\quad A:\ P(S_{done}).
\]

\paragraph{Barrier}
若 $N$ 个线程必须全部完成阶段 1 后才能进入阶段 2，核心是不变量：
\[
\text{phase2 begins}\Rightarrow count=N.
\]
简单一次性 barrier 可用受保护计数器 + gate 信号量构造；可复用 barrier 则需要防止下一轮线程混入上一轮，属于更高阶变式。

\section{从 PV 到管程：wait 为什么必须“释放并睡眠”}

\subsection{问题不是“wait 会不会释放锁”，而是不能丢失唤醒}

设线程 $T_1$ 检查：
\[
count=0
\]
准备等待。如果“释放锁”和“真正进入等待队列”之间存在可被其他线程插入的空窗：

\begin{codeblock}
T1: check count == 0
                 <- context switch
T2: count = 1
T2: signal()
T1: sleep
\end{codeblock}

那么通知发生时 T1 还没睡，之后却永久睡下，形成\kw{lost wakeup}。

因此条件等待必须具有语义：
\[
\boxed{\text{atomically release mutex and block}}
\]
被唤醒后，又必须：
\[
\boxed{\text{reacquire mutex before wait returns}}
\]
这样线程重新检查共享状态时仍处于受保护环境。

\subsection{标准思维模板}

\begin{codeblock}
lock(m);
while (!predicate)
    wait(cv, m);   // 原子释放 m + 阻塞；返回前重新获得 m
use_state();
unlock(m);
\end{codeblock}

修改方：
\begin{codeblock}
lock(m);
change_state();
signal(cv);
unlock(m);
\end{codeblock}

\begin{warn}[两个必须纠正的直觉]
\begin{itemize}
\item \code{signal()} 并不等于“条件一定成立且被唤醒者立即运行”；它通常只是使等待者有机会变为可运行状态。
\item 条件变量本身不是条件；真正的条件是共享状态上的布尔谓词，所以等待者醒来后必须用 \code{while} 重检。
\end{itemize}
\end{warn}

\subsection{条件变量 \code{signal()} 与信号量 V 不能互换}

两者都可能唤醒等待者，但拥有的状态不同：信号量把“许可数量”保存在自身计数中；条件变量只保存等待关系，真正的条件由受 mutex 保护的共享状态拥有。

\begingroup
\small
\renewcommand{\arraystretch}{1.25}
\setlength{\tabcolsep}{5pt}
\begin{tabularx}{\linewidth}{>{\bfseries}p{0.18\linewidth} X X X}
\toprule
操作 & 修改什么 & 当前无人等待 & 被唤醒者接下来做什么 \\
\midrule
信号量 \code{V(S)} & 增加一份许可；若有等待者则可交付许可并唤醒一个 & 许可被计数保存，之后的 \code{P(S)} 可以消费 & 获得许可后继续，具体调度时点另论 \\
条件变量 \code{signal(cv)} & 不替调用者修改共享谓词，只通知等待队列 & 通知通常不被保存 & 先成为可运行候选，并在 \code{wait} 返回前重新获得 mutex、重检谓词 \\
\bottomrule
\end{tabularx}
\endgroup

\begin{method}[稳定推理顺序]
写管程时先由调用者在锁内修改共享状态，再判断是否需要 \code{signal/broadcast}。写信号量时，许可计数本身就是同步状态的一部分。只因为两者都出现“唤醒”就互换，会同时破坏资源计数与条件谓词的 Owner。
\end{method}

\subsection{管程内的三类等待关系：入口、条件与紧急队列不能混名}

在经典 Hoare 管程实现中，等待实体至少分三类：
\begin{itemize}
\item \kw{Entry Queue}：尚未取得管程互斥使用权的外部调用者；
\item \kw{Condition Queue $Q_x$}：已在管程中检查谓词 $x$ 不成立，执行 \code{x.wait} 后释放管程并等待条件变化的实体；
\item \kw{Urgent/Next Queue}：发出 \code{x.signal} 后把管程立即交给被唤醒者，signal caller 暂时等待取回管程的位置。
\end{itemize}
它们的 WaitFor Relation 不同：入口队列等“管程可进入”，条件队列等“共享谓词可能变化”，紧急队列等“被我唤醒的实体离开或再次等待”。把三者都叫“阻塞队列”会丢失唤醒对象与交权优先级。

Hoare 语义的经典握手可压成：
\[
\text{signal caller}\xrightarrow{\text{wake }Q_x}\text{urgent wait},
\qquad
\text{awakened waiter}\xrightarrow{\text{direct handoff}}\text{monitor owner}.
\]
因此 waiter 醒来后不重新与 Entry Queue 竞争；它在 signal 与恢复之间直接拥有管程，刚成立的谓词不会被第三者插入破坏。退出管程时通常优先把控制权还给 urgent queue，再开放 entry queue。这里描述的是 Hoare 教材实现，不是要求所有条件变量库都真的暴露三种 queue 对象。

\begin{eng}[Hoare 与 Mesa 语义边界]
不同管程模型对 \code{signal} 后谁立即获得管程有不同约定。现代常见 Mesa 风格让 signal caller 继续，等待者只变为以后重新竞争 mutex 的候选，因此必须用 \code{while} 重检；Hoare 风格把控制权直接交给被唤醒者。Hansen 的经典折中要求 signal 放在管程过程末尾，使 caller 不再需要继续占有管程。考题未声明具体语义时，保留“通知不等于条件永久成立、唤醒不等于立即运行”这条稳定结论。
\end{eng}

\subsection{管程是什么}

管程可以概括为：
\[
\boxed{
\text{Shared State}
+
\text{Operations}
+
\text{Implicit Mutual Exclusion}
+
\text{Condition Variables}
}
\]
它的核心价值是把“共享状态及其合法同步协议”封装在一个模块中，减少调用者随意破坏不变量的机会。

\section{死锁：从等待依赖到资源状态安全性}

\begin{corebox}[死锁分析核心控制流]
\[
\boxed{
\begin{aligned}
\text{Resource State} &\to \text{Wait Dependency} \to \text{Deadlock Conditions} \to \text{State Classification} \\
&\to \text{Prevention / Avoidance / Detection} \to \text{Recovery} \to \text{Cost}
\end{aligned}
}
\]
面对任何死锁题目，不再从盲目记忆算法步骤开始，而是固定问 5 个核心问题：
\begin{enumerate}
\item \kw{谁持有什么？}（$\text{Allocation}$）
\item \kw{谁还在等什么？}（$\text{Request} / \text{Need}$）
\item \kw{这个等待有没有形成不可消解的依赖？}（Wait Dependency / Cycle）
\item \kw{题目是在问 Prevention、Avoidance 还是 Detection？}
\item \kw{我现在掌握的是未来最大需求 $\text{Need}$，还是当前未满足请求 $\text{Request}$？}
\end{enumerate}
\end{corebox}

\subsection{死锁到底是什么}

死锁 (Deadlock) 是多个执行流形成无法自行解除的资源等待环。每个执行流都等待其他执行流持有的不可剥夺资源，因此全部阻塞，系统无法继续推进。

死锁不等于长时间等待 (Long wait)。长时间等待是暂时的，一旦持有资源的线程释放锁即可唤醒；而死锁的等待链形成了环且无法自动解除，属于永久性阻塞。

\begin{method}[死锁推理的前置假设：资源类型区分]
分析死锁时，不能默认系统资源都是同一种模型，必须区分：
\begin{itemize}
  \item \kw{Reusable Resource（可重用资源）}：如 CPU、内存页、物理设备、锁/Mutex。一次只能被一个进程使用，使用完毕后被释放重用，总数保持不变。经典银行家算法与死锁检测建立在此模型上。
  \item \kw{Consumable Resource（可消耗资源）}：如消息 (Message)、信号 (Signal)、事件许可 (Event Token)。由生产者创建，被消费者接收后即被消费消耗，数量动态变化。
\end{itemize}
\end{method}

\begingroup
\small
\renewcommand{\arraystretch}{1.25}
\setlength{\tabcolsep}{6pt}
\begin{tabularx}{\linewidth}{>{\bfseries}p{0.22\linewidth} X}
\toprule
异常并发现象 & 直接原因与系统表现\\
\midrule
死锁 (Deadlock) & 多个执行流因相互等待对方持有的许可而全部阻塞，系统整体与个体均无法推进\\
活锁 (Livelock) & 执行流未阻塞且持续改变状态/重试，但由于缺乏协调（如持续对称退让）导致整体无法推进\\
饥饿 (Starvation) & 系统整体在推进（无死锁），但特定执行流因调度策略不公长期拿不到资源\\
\bottomrule
\end{tabularx}
\endgroup

\subsection{等待图与资源分配图}

分析死锁依赖时，不能仅看代码表面，必须区分两类依赖图：

\begin{enumerate}
\item \kw{Wait-for Graph（等待图）}：节点仅包含进程 $P_i$。边 $P_i \to P_j$ 表示 $P_i$ 正在等待只能由 $P_j$ 释放的资源。
\item \kw{Resource-Allocation Graph（资源分配图）}：节点包含进程 $P_i$ 与资源类 $R_j$。请求边 $P_i \to R_j$ 表示 $P_i$ 正在申请资源 $R_j$；分配边 $R_j \to P_i$ 表示 $R_j$ 已被分配给 $P_i$。
\end{enumerate}

在不同资源实例模型下，图中是否存在“环 (Cycle)”的判定标准截然不同：

\begin{itemize}
\item \kw{单实例资源模型（每类资源仅 1 个实例）}：等待图或分配图中的环是死锁的充要条件：
\[
\boxed{\text{Single-Instance: } \text{Cycle} \iff \text{Deadlock}}
\]
\item \kw{多实例资源模型（每类资源有多个实例）}：环仅表示存在死锁风险（必要不充分条件）。环内进程可能等待环外持有同类资源且即将结束的进程释放资源，从而打破环。
\end{itemize}

因此，考场上必须牢记第一处判定边界：
\[
\boxed{\text{Cycle} \neq \text{Universal Deadlock Test}}
\]

\subsection{Coffman 四个必要条件}

死锁形成必须同时满足 Coffman 四个必要条件：
\[
\boxed{\text{Deadlock} \implies \text{ME} \land \text{HW} \land \text{NP} \land \text{CW}}
\]

\begin{enumerate}
\item \kw{互斥 (Mutual Exclusion, ME)}：资源一次只能被一个进程独占使用。
\item \kw{请求并保持 (Hold and Wait, HW)}：进程至少持有一份资源，又提出了新资源请求，同时不释放已有资源。
\item \kw{不可剥夺 (No Preemption, NP)}：资源只能由持有进程自愿释放，系统不能强行剥夺。
\item \kw{循环等待 (Circular Wait, CW)}：存在进程等待链 $\{P_0, P_1, \dots, P_n\}$，使得 $P_i$ 等待 $P_{i+1}$ 持有的资源。
\end{enumerate}

学习这四个条件不要死记硬背名字，而要回答：\kw{为什么缺了一条件，依赖环就无法闭合？}

例如对于不可剥夺 (NP)：如果资源可以被系统安全地剥夺并恢复现场，当 $P_1$ 请求 $R_2$ 受阻时，OS 可以剥夺 $P_1$ 已持有的 $R_1$ 并分配给其他进程，依赖链就被打破了！

需要特别注意准确的边界表述：
\[
\boxed{\text{仅竞争可安全抢占的资源} \implies \text{不会因这些资源形成经典资源死锁}}
\]
不能绝对化地写成“系统只要出现可抢占资源就不会死锁”，因为进程完全可能同时持有着其他不可抢占的资源。

\subsection{死锁处理的四种策略层级}

\begingroup
\small
\renewcommand{\arraystretch}{1.25}
\setlength{\tabcolsep}{6pt}
\begin{tabularx}{\linewidth}{>{\bfseries}p{0.25\linewidth} X}
\toprule
死锁处理策略 & 核心逻辑与适用场景\\
\midrule
忽略 (Ostrich / 鸵鸟策略) & 假装死锁不发生。适用于死锁发生概率极低、预防/检测开销过高时（如通用桌面 OS）\\
预防 (Prevention) & 在系统设计期破坏 Coffman 四条件之一，静态保证系统绝不发生死锁\\
规避 (Avoidance) & 在运行时动态评估资源请求，若分配后系统处于安全状态才批准（如 Banker's Algorithm）\\
检测与恢复 (Detection \& Recovery) & 允许死锁发生，后台定期运行检测算法，发现死锁后通过剥夺/撤销进程恢复\\
\bottomrule
\end{tabularx}
\endgroup

需要澄清一个工程误区：通用操作系统不运行全局 Banker 算法，不等于完全不做死锁工程防范：
\[
\boxed{\text{No universal Banker} \neq \text{No deadlock engineering}}
\]
例如 Linux 内核广泛采用固定的锁申请顺序 (Lock Ordering)，并使用 \code{lockdep} 工具在运行时对锁依赖环进行静态/动态检查。

\subsection{Prevention：破坏必要条件的机制与代价}

\begingroup
\small
\renewcommand{\arraystretch}{1.25}
\setlength{\tabcolsep}{6pt}
\begin{tabularx}{\linewidth}{>{\bfseries}p{0.22\linewidth} p{0.38\linewidth} X}
\toprule
破坏的条件 & 典型实现机制 & 核心工程代价\\
\midrule
Mutual Exclusion & 将独占资源虚拟化/共享化（如 SPOOLing） & 许多物理资源本质上无法共享\\
Hold and Wait & 静态一次性申请全部资源；或申请新资源前释放已有资源 & 资源利用率极低、等待时间显著增加\\
No Preemption & 申请受阻时释放已持有的资源，稍后重新申请 & 仅适用于易保存恢复现场的资源（内存/寄存器），有 Rollback 开销\\
Circular Wait & 给资源建立全局唯一编号，强制按递增顺序申请 & 降低编程灵活性，增加模块间耦合\\
\bottomrule
\end{tabularx}
\endgroup

\begin{method}[SPOOLing 概念精确说明]
SPOOLing 并没有把物理打印机本身变成共享资源。物理打印机依旧是独占设备；SPOOLing 是通过守护进程 (Spooler Daemon) 集中接管物理设备，使多个用户进程可以并发提交逻辑打印任务。
\end{method}

在真实工程中，\kw{全序资源申请 (Global Resource Ordering)} 是最广泛采用的死锁预防思想（例如 Linux 目录锁机制规定必须按 inode 地址或层级顺序加锁）。

\subsection{State Classification：Safe / Unsafe / Deadlocked 集合关系}

系统的资源分配状态可划分为三个包含关系的集合：

\[
\boxed{\text{Deadlocked} \subset \text{Unsafe} \subset \text{All Allocation States}}
\]

即：
\[
\boxed{\text{Deadlocked} \implies \text{Unsafe}} \quad \text{但} \quad \boxed{\text{Unsafe} \nRightarrow \text{Deadlocked}}
\]

\begin{corebox}[Unsafe 状态表示什么]
Unsafe State 不等于此刻系统已经陷入死锁！Unsafe 的真实含义是：\kw{从当前状态出发，系统无法保证无论未来进程如何提出请求，都一定存在一条能让所有进程顺利完成的安全序列}。

如果未来进程实际提出的资源请求少于其声明的最大需求 $\text{Need}$，处于 Unsafe 状态的系统完全可能安全运行完毕而不发生死锁。
\end{corebox}

\subsection{单类可重用资源的最小总量公式：先写适用条件}

若题目同时满足：只有一类同质、不可剥夺、可重用资源；共有 $n$ 个进程；进程 $P_i$ 的最大需求为 $m_i\ge 1$；进程获得全部所需资源后能够完成并一次释放；题目要求的是“无论怎样逐步占有，都保证至少有一个进程可继续”，则最小资源总量为
\[
\boxed{R_{\min}=1+\sum_{i=1}^{n}(m_i-1)=\sum_{i=1}^{n}m_i-n+1.}
\]

最坏状态是每个进程都持有 $m_i-1$ 份并各差一份。再多一份资源，至少可使一个进程达到最大需求并完成；它释放后形成连锁推进。

\begin{warn}[这个公式不是通用 Banker]
只要出现多类资源向量、当前 Allocation/Need 矩阵、可部分释放、资源可剥夺，或题目问某个\kw{给定状态}是否安全，就必须回到安全性算法，不能套标量公式。公式的价值是压缩一个严格受限的单资源模型，不是绕过状态建模。
\end{warn}

\subsection{Avoidance 与银行家算法 (Banker's Algorithm)}

银行家算法先临时执行一次资源分配，再检查新状态是否仍存在安全序列。

\paragraph{核心数据结构}
设系统有 $n$ 个进程，$m$ 类资源：
\begin{itemize}
\item $\text{Available}[j]$：资源 $R_j$ 当前未被分配的可用实例数。
\item $\text{Max}[i, j]$：进程 $P_i$ 对资源 $R_j$ 的最大需求数。
\item $\text{Allocation}[i, j]$：进程 $P_i$ 当前已分配到的资源 $R_j$ 实例数。
\item $\text{Need}[i, j] = \text{Max}[i, j] - \text{Allocation}[i, j]$：进程 $P_i$ 未来还需要资源 $R_j$ 的最大数量。
\end{itemize}

\paragraph{资源请求 $\text{Request}_i$ 的三关校验}
当进程 $P_i$ 发出资源请求 $\text{Request}_i$ 时，必须依次通过三关：

\begin{enumerate}
\item \kw{第一关：合法性校验}
\[
\text{Request}_i \le \text{Need}_i
\]
若违反，说明请求超过了其声明的最大需求，抛出非法请求错误。

\item \kw{第二关：充沛性校验}
\[
\text{Request}_i \le \text{Available}
\]
若违反，说明“当前可用资源不足（$\boxed{\text{Insufficient Now}}$）”，$P_i$ 必须阻塞等待。注意：此时不是 Unsafe，只是暂时无法满足。

\item \kw{第三关：试探分配与安全性测试}
系统假设批准请求，试探性更新状态：
\[
\begin{aligned}
\text{Available}' &= \text{Available} - \text{Request}_i \\
\text{Allocation}_i' &= \text{Allocation}_i + \text{Request}_i \\
\text{Need}_i' &= \text{Need}_i - \text{Request}_i
\end{aligned}
\]
随后调用安全性算法 $\boxed{\text{SafetyTest}(\text{State}')}$：
\begin{itemize}
\item 若 $\text{State}'$ 安全 $\implies$ \kw{Commit}（正式批准分配）；
\item 若 $\text{State}'$ 不安全 $\implies$ \kw{Rollback}（撤销试探分配，$P_i$ 恢复阻塞等待）。
\end{itemize}
\end{enumerate}

\[
\boxed{
\text{Tentative Transition}
\rightarrow
\text{Safety Check}
\rightarrow
\text{Commit / Rollback}
}
\]

\subsection{安全算法 (Safety Algorithm) 极简理解}

安全算法通过模拟“进程按某种顺序依次获得所需全部资源并运行结束归还资源”来寻找安全序列：

\begin{enumerate}
\item 初始化：$\text{Work} = \text{Available}$，$\text{Finish}[i] = \text{false} \ (i=1,\dots,n)$。
\item 寻找满足条件的进程 $P_i$：
\[
\text{Finish}[i] == \text{false} \quad \land \quad \text{Need}_i \le \text{Work}
\]
\item 若找到，模拟 $P_i$ 拿到 $\text{Need}_i$ 后运行完成，归还其持有的全部资源：
\[
\text{Work} \leftarrow \text{Work} + \text{Allocation}_i, \quad \text{Finish}[i] = \text{true}
\]
返回步骤 2。
\item 若所有进程 $\text{Finish}[i] == \text{true}$，则系统处于安全状态；若最终仍有进程 $\text{Finish}[i] == \text{false}$，则处于不安全状态。
\end{enumerate}

\begin{corebox}[为什么 Work 递增的是 Allocation 而不是 Max？]
学生最容易混淆为什么 $\text{Work}' = \text{Work} + \text{Allocation}_i$。

因为 $P_i$ 运行需要从 $\text{Work}$ 中拿走它还缺少的 $\text{Need}_i$。$P_i$ 执行完成后归还的是它拿走的 $\text{Need}_i$ 以及原先已持有的 $\text{Allocation}_i$（共归还 $\text{Need}_i + \text{Allocation}_i$）。

系统之前已经扣除了 $\text{Need}_i$，因此归还后系统的净资源增加量恰好就是 $P_i$ 原先占有的部分：
\[
\boxed{\text{Work}' = (\text{Work} - \text{Need}_i) + (\text{Need}_i + \text{Allocation}_i) = \text{Work} + \text{Allocation}_i}
\]
\end{corebox}

\begin{method}[银行家手算的三个自检不变量]
只给总量时，先用
\[
\boxed{\text{Available}=\text{Total}-\sum_i\text{Allocation}_i}
\]
核对 Work 初值；每完成一个安全序列候选，只执行
\[
\text{Work}\leftarrow\text{Work}+\text{Allocation}_i,
\]
因为候选还缺的 Need 是模拟中先借后还的中间量。Work 只增不减，所以当前满足 $\text{Need}_i\le\text{Work}$ 的进程以后不会失去可完成性；选任意一个可推进者不会漏掉安全序列，跑不动即不存在安全序列。若所有进程都完成，末态必须回到 $\text{Total}$；否则优先检查 Allocation 求和、Need 减法或向量比较。

安全算法每轮扫描 $n$ 个进程并比较 $m$ 类资源，朴素实现上界为 $O(n^2m)$；它需要进程声明 Max，且动态进程加入、Max 不可预知、资源类型变化都会削弱适用性。它按最坏需求做保证，可能保守拒绝实际不会死锁的请求，这正是 Avoidance 与事后 Detection 的资源利用率权衡。
\end{method}

\subsection{Avoidance 与 Detection：目标状态与输入语义不同}

死锁规避 (Banker) 与 死锁检测 (Detection) 在算法结构上相似，但语义与触发时机有本质不同。考场上必须对比澄清：

\begingroup
\small
\renewcommand{\arraystretch}{1.25}
\setlength{\tabcolsep}{6pt}
\begin{tabularx}{\linewidth}{>{\bfseries}p{0.22\linewidth} X X}
\toprule
对比维度 & 死锁规避 (Avoidance / Banker) & 死锁检测 (Deadlock Detection)\\
\midrule
触发时机 & 资源分配之前（事前试探） & 允许分配之后（事后定期/阻塞时检测）\\
核心目的 & 保证系统绝不进入 Unsafe 状态 & 判断系统当前是否已经陷入死锁\\
关键输入矩阵 & $\text{Need}$（未来最大需求） & $\text{Request}$（当前实际正在等待的请求）\\
矩阵语义 & 进程未来最多还要再申请多少 & 进程此刻卡住正在等多少\\
推演成功含义 & 存在安全序列，请求可批准 & 可消去所有进程，当前没有死锁\\
推演失败含义 & Unsafe，拒绝当前分配试探 & 剩余无法消去的进程处于死锁状态\\
一句话总结 &
$\boxed{\begin{gathered}\text{Banker: Can everybody}\\\text{still finish?}\end{gathered}}$ &
$\boxed{\begin{gathered}\text{Detection: Who is}\\\text{already unable to finish?}\end{gathered}}$\\
\bottomrule
\end{tabularx}
\endgroup

\subsection{Deadlock Detection（死锁检测算法）}

死锁检测算法不预测未来最大需求 $\text{Need}$，只关注\kw{此刻卡住的实际请求 $\text{Request}_i$ 能否被满足}。

\paragraph{多实例资源检测（消去法）}
\begin{enumerate}
\item 初始化：$\text{Work} = \text{Available}$。对于 $\text{Allocation}_i \neq 0$ 的进程，$\text{Finish}[i] = \text{false}$；若 $\text{Allocation}_i == 0$，则 $\text{Finish}[i] = \text{true}$。
\item 寻找进程 $P_i$ 满足：
\[
\text{Finish}[i] == \text{false} \quad \land \quad \text{Request}_i \le \text{Work}
\]
\item 若找到，说明 $P_i$ 当前请求可被满足，模拟其获得资源运行结束并释放已占有的资源：
\[
\text{Work} \leftarrow \text{Work} + \text{Allocation}_i, \quad \text{Finish}[i] = \text{true}
\]
重复步骤 2。
\item 算法终止时，若存在 $\text{Finish}[i] == \text{false}$ 的进程，则这些进程处于死锁状态。
\end{enumerate}

消去法之所以不需要 $\text{Max}$ 矩阵，就是因为它只问“现在卡住的请求能否解决”，一旦解决即可假定其能归还资源。

\subsection{Recovery：死锁恢复与代价评估}

当检测算法确认存在死锁后，系统必须执行恢复。恢复途径分为进程终止与资源剥夺：

\begin{enumerate}
\item \kw{进程终止 (Process Termination)}：
\begin{itemize}
\item 一次性终止所有死锁进程（简单但代价极大）；
\item 逐个终止进程，每终止一个重新运行检测算法，直到死锁环消除。
\end{itemize}
\item \kw{资源抢占 (Resource Preemption)}：
\begin{itemize}
\item 涉及三个核心决策：选择牺牲者 (Victim selection)、回滚 (Rollback) 与 避免饥饿 (Prevent starvation)。
\end{itemize}
\end{enumerate}

恢复决策统一遵循代价最小化函数：
\[
\boxed{\text{Recovery Cost} = \text{LostWork} + \text{RollbackCost} + \text{ResourceCost} + \text{StarvationRisk}}
\]

选择牺牲者进程 (Victim) 时综合评估 6 个维度：进程优先级、已运行时间与已完成工作量、尚需多少时间完成、当前持有多少资源、终止可释放多少资源、已被回滚/剥夺的次数。

\subsection{检测不是“利用率低就判死锁”：触发时机与恢复检查点}

检测可以在每次资源请求后做，发现快但成本高；也可以定时做，周期越长死锁占用资源和卷入进程越多；还可以在“利用率明显下降且就绪队列并不长”等征兆出现时按需做。利用率低本身不能证明死锁，I/O 密集型负载同样会造成低 CPU 使用，最终仍须运行图化简/矩阵检测。

逐个终止死锁进程时，每终止一个都要重新检测，因为一次释放可能打破整个等待环；资源抢占则必须恢复被抢进程到一致状态。没有检查点时只能完全重做，存在检查点时可回滚到最近安全快照；检查点越密，恢复损失越小但正常保存成本越高。牺牲者选择还应记录已被选中过的次数，否则同一低优先级进程会反复被抢而形成恢复饥饿。

\section{专题总结：把 PV 从代码模板提升为状态约束与语义建模}

\begin{corebox}[最终核心模型]
\[
\boxed{
\begin{aligned}
\text{Flows}&\to\text{Shared State}\to\text{Interleavings}\to\text{Invariant}\to\text{Constraints}\\
&\to\text{Semaphore/PV}\to\text{Progress}\to\text{Cost}
\end{aligned}}
\]
\end{corebox}

不要看到生产者--消费者就机械写 \code{empty/full/mutex}。按以下顺序建模：
\begin{enumerate}
\item 先找到\kw{谁在等待什么事实}；
\item 再找到\kw{谁能够制造这个事实}；
\item 用信号量把事实变成可被原子消费/生产的许可；
\item 用互斥保护共享状态的不变量；
\item 最后证明所有允许交错下系统既不会进入坏状态，也不会因为等待环而无法推进。
\end{enumerate}

\begin{method}[与进程调度的接口]
P 操作可能使执行流从 Running 进入 Blocked，V 操作可能使等待者从 Blocked 进入 Ready；Ready$\to$Running 仍由调度器决定。PV 只表达许可和等待，不直接分配 CPU。
\end{method}

\section*{边界检查：四句不能绝对化}
\addcontentsline{toc}{section}{边界检查：四句不能绝对化}

\begin{warn}[四句不应绝对化的话]
\begin{enumerate}
\item “所有 PV 题都必须先同步 P 后互斥 P” —— 标准有界缓冲区正确，但更一般原则是避免持锁等待依赖环。
\item “V 的顺序永远不重要” —— 简单模板里可能不影响 Safety，但会影响唤醒竞争、性能，复杂协议还可能影响语义。
\item “二元信号量就是 mutex” —— 考试可用于互斥，但严格语义中 mutex 有所有权，semaphore 表示许可。
\item “某个哲学家解法一定无饥饿” —— 无死锁不自动等于公平；饥饿还取决于调度和唤醒策略。
\end{enumerate}
\end{warn}

\end{document}
